% Table 2 (3.2 Unimodal performance)
% 자동 생성: tools/make_paper_tables.py — 수치를 여기서 고치지 말고 스크립트를 다시 돌릴 것.
% 단독 문서다(캡처용). Overleaf 새 프로젝트에서 pdfLaTeX 로 컴파일한다.
% 기존 main.tex 에 넣을 때는 이 파일이 아니라 overleaf_tables.tex 를 쓴다.
\documentclass[border=8pt]{standalone}
\usepackage[T1]{fontenc}
\usepackage{newtxtext,newtxmath}
\usepackage{booktabs,multirow,array,pifont}
% --- 캡처할 때 제목·주석을 빼려면 아래 두 줄을 \...false 로 바꾼다 ---
\newif\ifshowtitle \showtitletrue
\newif\ifshownotes \shownotestrue
\newsavebox\tablebox
\begin{document}
\sbox\tablebox{%
\begin{tabular}{@{}llcc@{}}
\toprule
Modality & Encoder & OS C-index & PFS C-index \\
\midrule
Clinical variables & MLP & \textbf{0.639 $\pm$ 0.044} & \textbf{0.622 $\pm$ 0.037} \\
\addlinespace
\multirow{4}{*}{Radiology report} & TF-IDF & 0.627 $\pm$ 0.047 & 0.609 $\pm$ 0.038 \\
 & mBERT & 0.640 $\pm$ 0.024 & 0.567 $\pm$ 0.056 \\
 & KM-BERT & 0.607 $\pm$ 0.068 & 0.614 $\pm$ 0.021 \\
 & RadBERT & \textbf{0.668 $\pm$ 0.037} & \textbf{0.635 $\pm$ 0.042} \\
\addlinespace
\multirow{3}{*}{PET-CT MIP} & ResNet-18 (ImageNet) & 0.633 $\pm$ 0.042 & 0.572 $\pm$ 0.018 \\
 & ResNet-50 (RadImageNet) & 0.616 $\pm$ 0.050 & 0.556 $\pm$ 0.046 \\
 & Simple CNN & \textbf{0.657 $\pm$ 0.048} & \textbf{0.615 $\pm$ 0.057} \\
\bottomrule
\end{tabular}%
}%
\begin{minipage}{\wd\tablebox}
\ifshowtitle\noindent \textbf{Table 2.} Prognostic Performance of Unimodal Models\par\vspace{4pt}\fi
\noindent\usebox\tablebox
\ifshownotes\par\vspace{3pt}\footnotesize\noindent Mean $\pm$ SD of the C-index across the five test folds (n = 238). Bold, encoder used in the multimodal models. CNN = convolutional neural network, KM-BERT = Korean medical BERT, mBERT = multilingual BERT, MIP = maximum intensity projection, MLP = multilayer perceptron, OS = overall survival, PFS = progression-free survival, SD = standard deviation, TF-IDF = term frequency--inverse document frequency\par\fi
\end{minipage}
\end{document}
